\documentclass[uplatex,a4paper,11pt]{jsarticle}
\usepackage[margin=22mm]{geometry}
\usepackage[dvipdfmx]{hyperref,xcolor}
\usepackage{fvextra,enumitem,longtable,booktabs,amssymb}
\usepackage{listings,jvlisting}
\hypersetup{colorlinks=true,linkcolor=blue!55!black,urlcolor=blue!65!black,
  pdftitle={P1 TeXNote・LilyPondNote認証サーバー構築手順}}
\DefineVerbatimEnvironment{command}{Verbatim}{fontsize=\small,frame=single,
  rulecolor=\color{gray!55},breaklines=true,breakanywhere=true}
\setlist{nosep}
\lstset{language=Python,basicstyle=\ttfamily\scriptsize,numbers=left,
  numberstyle=\tiny,frame=single,breaklines=true,columns=fullflexible,
  keepspaces=true,showstringspaces=false,tabsize=4,captionpos=b}

\title{P1 TeXNote・LilyPondNote\\認証サーバー構築手順\\
\large DB・venv・ソースの対応が分かる完成構成版}
\author{}
\date{2026年8月31日}

\begin{document}
\maketitle

\begin{abstract}
本書は、P1にTeXNote認証サービスとLilyPondNote認証サービスを最初から完成形で構築する
手順である。LilyPondNoteの\texttt{main.py}にはPDF生成用\texttt{/compile}と移調
テキスト生成用\texttt{/transpose}と、Free・Standard・Proの契約プラン制があらかじめ
実装されている。本書ではこの完成構成を新規に構築する。二つのDB、二つのvenv、requirements、配置する
ソースの対応関係を先に説明してから、実際の構築手順を示す。
\end{abstract}

\tableofcontents
\newpage

\section{最初に理解する完成構成}

\subsection{P1には二つの独立した認証サービスがある}

\begin{command}
P0
 |-- tex.example.com      -> P1 TeX認証API       TCP 8000
 |                            DB: /opt/texnote/P1/tex_service.db
 |                            venv: /opt/texnote/P1/.venv
 |
 `-- lilypond.example.com -> P1 LilyPond認証API  TCP 8001
                              DB: /opt/lilypondnote/P1/lilypondnote.db
                              venv: /opt/lilypondnote/P1/.venv
                              |-- /compile   -> P2:8001 PDF生成
                              `-- /transpose -> P2:8002 移調テキスト生成
\end{command}

P1はP0の機能を持たない。Nginx、TLS証明書、インターネットからのポート転送をP1へ
設定しない。

\subsection{何と何が別なのか}

\begin{longtable}{@{}p{40mm}p{51mm}p{55mm}@{}}
\toprule
項目 & TeXNote & LilyPondNote\\
\midrule
OSユーザー & \texttt{texnote-p1} & \texttt{lilypondnote-p1}\\
Python venv & \path{/opt/texnote/P1/.venv} & \path{/opt/lilypondnote/P1/.venv}\\
SQLite DB & \texttt{tex\_service.db} & \texttt{lilypondnote.db}\\
JWT署名鍵 & TeX専用 & LilyPond専用\\
P2内部トークン & TeX専用 & LilyPondの2処理で共用\\
P1待受 & TCP 8000 & TCP 8001\\
\bottomrule
\end{longtable}

したがって、\textbf{TeXNoteとLilyPondNoteは別DB}である。同じ人が両方を利用する場合は、
両DBへそれぞれ利用者を登録する。

一方、LilyPondのPDF生成と移調テキスト生成は別DBではない。\textbf{一つの
\texttt{lilypondnote.db}、一つの利用者アカウント、一つのJWT、一つの契約プランと月間上限を共有}
する。どちらを利用したかは、同じ\texttt{usage}テーブルの\texttt{service}列へ
\texttt{typeset}または\texttt{transpose}と記録する。

LilyPondNoteアプリの移調設定には「ローカル」と「リモート」がある。ローカル移調は端末内で
完結してP1へ要求を送らないため、利用回数へ加算されない。リモート移調だけが
\texttt{/transpose}を通り、成功・失敗を含む1要求として月間合計と移調内訳へ記録される。

\section{二つのDB構造}

\subsection{TeXNote DB}

保存先は\path{/opt/texnote/P1/tex_service.db}である。

\begin{longtable}{@{}p{28mm}p{46mm}p{72mm}@{}}
\toprule
テーブル & 主な列 & 用途\\
\midrule
\texttt{users} & id, email, password\_hash, status, plan & TeX利用者、状態、契約プラン\\
\texttt{usage} & user\_id, created\_at, compile\_seconds, input\_bytes, success & TeXコンパイル履歴\\
\texttt{api\_keys} & user\_id, key\_hash, revoked\_at & 将来のAPIキー管理用\\
\bottomrule
\end{longtable}

TeX側の\texttt{usage}にLilyPondの履歴は入らない。

\subsection{LilyPondNote DB}

保存先は\path{/opt/lilypondnote/P1/lilypondnote.db}である。

\begin{longtable}{@{}p{28mm}p{49mm}p{69mm}@{}}
\toprule
テーブル & 列 & 用途\\
\midrule
\texttt{users} & id, email, password\_hash, status, plan & LilyPond利用者、状態、契約プラン\\
\texttt{usage} & id, user\_id, created\_at, compile\_seconds, input\_bytes, success, service & PDF生成・移調の共通履歴\\
\bottomrule
\end{longtable}

契約プランと月間上限は\texttt{plans.py}の一箇所で次のように定義する。PDF生成とリモート
移調は、この同じ上限を消費する。

\begin{longtable}{@{}p{45mm}p{55mm}@{}}
\toprule
プラン & 月間要求回数\\
\midrule
\texttt{free} & 50回\\
\texttt{standard} & 1,000回\\
\texttt{pro} & 5,000回\\
\bottomrule
\end{longtable}

\texttt{usage.service}の値は次の二つだけである。

\begin{longtable}{@{}p{35mm}p{105mm}@{}}
\toprule
値 & 意味\\
\midrule
\texttt{typeset} & \texttt{/compile}を使い、P2のLilyPondでPDFを生成した\\
\texttt{transpose} & \texttt{/transpose}を使い、P2で移調後テキストを生成した。PDFは生成しない\\
\bottomrule
\end{longtable}

新しいLilyPondNote DBを完成形で作成するソースをここに示す。

\lstinputlisting[caption={LilyPondNote DB作成・接続 database.py}]{/Users/mat/Documents/Git/LilyPondNote/Server/P1/database.py}

\section{OSパッケージとPythonパッケージの違い}

\subsection{aptで入れるもの}

aptはOS全体で使うPython本体、venv作成機能、SQLite確認コマンド、UFWを導入する。

\begin{command}
sudo apt update
sudo apt full-upgrade
sudo apt install python3 python3-venv python3-pip sqlite3 ufw
\end{command}

P1ではTeXやLilyPondを実行しないため、TeX LiveとLilyPond本体は導入しない。それらはP2へ
導入する。

\subsection{pipとrequirementsの役割}

pipは各venvの中だけへPythonライブラリを導入する。本書ではvenvを混用しないため、
requirementsも二つの明示的なファイルに分ける。さらに、各requirementsは、そのサービスの
ソースと同じリポジトリで管理する。したがってTeX用ファイルをLilyPondNoteリポジトリへ
置いてはならない。

\begin{longtable}{@{}p{48mm}p{48mm}p{48mm}@{}}
\toprule
requirements & 導入先venv & 用途\\
\midrule
TeXNoteリポジトリ\\\texttt{Server/P1/}\\\texttt{requirements-texnote.txt} & \path{/opt/texnote/P1/.venv} & TeX認証・JWT・P2中継\\
LilyPondNoteリポジトリ\\\texttt{Server/P1/}\\\texttt{requirements-}\\\texttt{lilypondnote.txt} & \path{/opt/lilypondnote/P1/.venv} & LilyPond認証・PDF中継・移調中継\\
\bottomrule
\end{longtable}

現在は両サービスが同じWeb技術を使うため内容は同じだが、venvは独立している。将来片方の
バージョンだけを変更できるよう、ファイルも分けて管理する。

\begin{longtable}{@{}p{35mm}p{105mm}@{}}
\toprule
パッケージ & 必要な理由\\
\midrule
FastAPI & HTTP APIと入力検査\\
Uvicorn & FastAPIを常時待受させるASGIサーバー\\
PyJWT & JWTの発行と検証\\
argon2-cffi & パスワードのArgon2ハッシュと照合\\
requests & P1からP2へ内部HTTP要求を送る\\
\bottomrule
\end{longtable}

標準ライブラリの\texttt{sqlite3}、\texttt{os}、\texttt{time}等をpipで入れてはならない。
multipart/form-dataを使わないため\texttt{python-multipart}も不要である。現在のソースに
対して、次のrequirementsが必要十分である。

\lstinputlisting[language={},caption={TeXNote P1 requirements}]{/Users/mat/Documents/Git/TeXNote/Server/P1/requirements-texnote.txt}

\lstinputlisting[language={},caption={LilyPondNote P1 requirements}]{/Users/mat/Documents/Git/LilyPondNote/Server/P1/requirements-lilypondnote.txt}

\section{専用ユーザーと配置先を作る}

\begin{command}
sudo useradd --system --home /opt/texnote/P1 \
  --shell /usr/sbin/nologin texnote-p1
sudo useradd --system --home /opt/lilypondnote/P1 \
  --shell /usr/sbin/nologin lilypondnote-p1
sudo install -d -o texnote-p1 -g texnote-p1 /opt/texnote/P1
sudo install -d -o lilypondnote-p1 -g lilypondnote-p1 /opt/lilypondnote/P1
\end{command}

\section{TeXNote認証サービスを構築する}

\subsection{必要なソースだけを配置する}

運転に必要な\texttt{main.py}、\texttt{auth.py}、\texttt{database.py}、
\texttt{plans.py}と、初期構築に必要な\texttt{tools/init\_db.py}、
\texttt{tools/useradd.py}、日常運用で使う\texttt{tools/check02.py}、
\texttt{tools/change\_plan.py}、\texttt{tools/reset\_password.py}を配置する。利用者の登録・確認・契約変更方法は、文書末尾の
「日常運用手続き」でまとめて説明する。

\begin{command}
# 管理用MacのTeXNoteリポジトリ直下で実行
scp Server/P1/progs/main.py Server/P1/progs/auth.py \
  Server/P1/progs/database.py Server/P1/progs/plans.py \
  mat@192.168.30.10:/tmp/
scp Server/P1/progs/tools/init_db.py Server/P1/progs/tools/useradd.py \
  Server/P1/progs/tools/check02.py Server/P1/progs/tools/change_plan.py \
  Server/P1/progs/tools/reset_password.py \
  mat@192.168.30.10:/tmp/

# P1で実行
sudo install -o texnote-p1 -g texnote-p1 -m 640 \
  /tmp/main.py /tmp/auth.py /tmp/database.py /tmp/plans.py \
  /opt/texnote/P1/
sudo install -d -o texnote-p1 -g texnote-p1 /opt/texnote/P1/tools
sudo install -o texnote-p1 -g texnote-p1 -m 640 \
  /tmp/init_db.py /tmp/useradd.py /tmp/check02.py /tmp/change_plan.py \
  /tmp/reset_password.py \
  /opt/texnote/P1/tools/
sudo -u texnote-p1 touch /opt/texnote/P1/tools/__init__.py
\end{command}

\subsection{TeX用venvを作る}

\begin{command}
# 管理用MacのTeXNoteリポジトリ直下で実行
scp Server/P1/requirements-texnote.txt mat@192.168.30.10:/tmp/

# P1で実行
sudo -u texnote-p1 python3 -m venv /opt/texnote/P1/.venv
sudo -u texnote-p1 /opt/texnote/P1/.venv/bin/pip install \
  --requirement /tmp/requirements-texnote.txt
sudo -u texnote-p1 /opt/texnote/P1/.venv/bin/python \
  -m py_compile /opt/texnote/P1/*.py /opt/texnote/P1/tools/*.py
sudo -u texnote-p1 sh -c \
  'cd /opt/texnote/P1 && .venv/bin/python -m tools.init_db'
\end{command}

利用者は初期構築中に仮登録せず、疎通試験の直前に「日常運用手続き」の登録手順で作成する。

\subsection{TeX側の必要ソース}

\lstinputlisting[caption={TeX認証・P2中継 main.py}]{/Users/mat/Documents/Git/TeXNote/Server/P1/progs/main.py}
\lstinputlisting[caption={TeXパスワード認証 auth.py}]{/Users/mat/Documents/Git/TeXNote/Server/P1/progs/auth.py}
\lstinputlisting[caption={TeX DB接続 database.py}]{/Users/mat/Documents/Git/TeXNote/Server/P1/progs/database.py}
\lstinputlisting[caption={TeXプラン定義 plans.py}]{/Users/mat/Documents/Git/TeXNote/Server/P1/progs/plans.py}
\lstinputlisting[caption={TeX DB初期化 tools/init\_db.py}]{/Users/mat/Documents/Git/TeXNote/Server/P1/progs/tools/init_db.py}

\section{LilyPondNote認証サービスを構築する}

\subsection{完成済みソースを最初から配置する}

ここで配置する\texttt{main.py}には、すでに次が実装されている。

\begin{itemize}
  \item \texttt{POST /login}：パスワード認証とJWT発行
  \item \texttt{GET /me}：契約プラン、上限、合計利用回数と\texttt{typeset}/\texttt{transpose}内訳
  \item \texttt{POST /compile}：P2 PDF APIへ中継し\texttt{typeset}履歴を記録
  \item \texttt{POST /transpose}：P2移調APIへ中継し\texttt{transpose}履歴を記録
  \item Free・Standard・Proの月間上限を、PDF生成とリモート移調の合計へ適用
\end{itemize}

ログインJSONには、次のサービス識別を必ず含める。

\begin{command}
TeXNote:       "service": "texnote"
LilyPondNote:  "service": "lilypondnote"
\end{command}

発行するJWTの\texttt{aud}にも同じサービス名を入れ、以後の\texttt{/me}、
\texttt{/compile}、\texttt{/transpose}で対象サービスの\texttt{aud}だけを受理する。
さらに\texttt{/me}は\texttt{service}を返す。この三段階により、P0のポート番号を
取り違えたクライアントが別サービスのDBでログインすることを防止する。
ここで掲載・配置する\texttt{main.py}が、認証、サービス識別、PDF生成中継、移調中継を
すべて備えた初期構築用の完成ソースである。

\begin{command}
# 管理用MacのLilyPondNoteリポジトリ直下で実行
scp Server/P1/main.py Server/P1/database.py Server/P1/plans.py \
  Server/P1/useradd.py Server/P1/change_plan.py Server/P1/reset_password.py \
  Server/P1/requirements-lilypondnote.txt \
  mat@192.168.30.10:/tmp/

# P1で実行
sudo install -o lilypondnote-p1 -g lilypondnote-p1 -m 640 \
  /tmp/main.py /tmp/database.py /tmp/plans.py /tmp/useradd.py \
  /tmp/change_plan.py /tmp/reset_password.py \
  /opt/lilypondnote/P1/
\end{command}

\subsection{LilyPond用venvを作る}

\begin{command}
sudo -u lilypondnote-p1 python3 -m venv /opt/lilypondnote/P1/.venv
sudo -u lilypondnote-p1 /opt/lilypondnote/P1/.venv/bin/pip install \
  --requirement /tmp/requirements-lilypondnote.txt
sudo -u lilypondnote-p1 /opt/lilypondnote/P1/.venv/bin/python \
  -m py_compile /opt/lilypondnote/P1/*.py
\end{command}

利用者を登録せず、まずDBと二つのテーブルだけを完成形で作成する。

\begin{command}
sudo -u lilypondnote-p1 sh -c \
  'cd /opt/lilypondnote/P1 && .venv/bin/python -c "from database import initialize; initialize()"'
sudo -u lilypondnote-p1 sqlite3 /opt/lilypondnote/P1/lilypondnote.db \
  '.schema'
\end{command}

\subsection{LilyPond側の必要ソース}

運転、初期登録、日常の契約変更に使う完成ソースを次に示す。

\lstinputlisting[caption={完成済みLilyPond認証・PDF/移調中継 main.py}]{/Users/mat/Documents/Git/LilyPondNote/Server/P1/main.py}
\lstinputlisting[caption={LilyPond DB作成・接続 database.py}]{/Users/mat/Documents/Git/LilyPondNote/Server/P1/database.py}
\lstinputlisting[caption={LilyPondプラン定義 plans.py}]{/Users/mat/Documents/Git/LilyPondNote/Server/P1/plans.py}
\lstinputlisting[caption={LilyPond利用者登録 useradd.py}]{/Users/mat/Documents/Git/LilyPondNote/Server/P1/useradd.py}
\lstinputlisting[caption={LilyPond契約プラン変更 change\_plan.py}]{/Users/mat/Documents/Git/LilyPondNote/Server/P1/change_plan.py}

\section{秘密値とP2接続先}

ここでは4個の異なるランダム値を作る。同じ値を使い回してはならない。それぞれの役割は
次のとおりである。

\begin{enumerate}
  \item \textbf{TeX用JWT署名鍵}：P1のTeX認証APIが、利用者へ発行するJWTへ署名し、
        後の要求でJWTが改変されていないことを確認する秘密鍵。
  \item \textbf{TeX用P2内部トークン}：P1のTeX認証APIからP2のTeX処理APIを呼ぶとき、
        正規のP1から来た内部要求であることをP2が確認するための共有トークン。
  \item \textbf{LilyPond用JWT署名鍵}：P1のLilyPond認証APIが、LilyPond利用者のJWTへ
        署名・検証するための秘密鍵。TeX用JWT署名鍵とは別物である。
  \item \textbf{LilyPond用P2内部トークン}：P1からP2のPDF生成APIと移調APIを呼ぶための
        共有トークン。この一つだけをLilyPondの二つの処理で共用する。
\end{enumerate}

管理用Macで\texttt{openssl rand -hex 32}を4回実行する。表示された順に上記1～4へ対応
させ、取り違えないよう一つずつ対象ファイルへ貼り付ける。

\begin{command}
openssl rand -hex 32
openssl rand -hex 32
openssl rand -hex 32
openssl rand -hex 32
\end{command}

Mac側には、あらかじめプレースホルダー入りの次の二ファイルを置いてある。

\begin{itemize}
  \item TeXNoteリポジトリ：\path{Server/P1/p1.env}
  \item LilyPondNoteリポジトリ：\path{Server/P1/p1.env}
\end{itemize}

TeXNote側\texttt{p1.env}の二つの\texttt{<...>}を上記1、2の値に、LilyPondNote側の
二つを上記3、4の値に置き換える。完成形は次の構造になる（下記の説明用記号そのものを
保存してはならない）。

TeXNoteリポジトリの\path{Server/P1/p1.env}：

\begin{command}
TEX_AUTH_SECRET_KEY=<TeX用JWT署名鍵>
TEXNOTE_API_TOKEN=<TeX用P2内部トークン>
TEXNOTE_P2_TYPESET_URL=http://192.168.40.10:8000/v1/typeset
\end{command}

LilyPondNoteリポジトリの\path{Server/P1/p1.env}：

\begin{command}
LILYPONDNOTE_AUTH_SECRET_KEY=<LilyPond用JWT署名鍵>
LILYPONDNOTE_API_TOKEN=<LilyPond用P2内部トークン>
LILYPONDNOTE_P2_TYPESET_URL=http://192.168.40.10:8001/v1/typeset
LILYPONDNOTE_P2_TRANSPOSE_URL=http://192.168.40.10:8002/v1/transpose
\end{command}

LilyPondのPDFと移調は同じP1認証サービスから送られるため、P2内部トークンも共用する。

編集後、Macから一度P1の\path{/tmp}へ転送する。同名ファイルなので一度に送らず、次の順で
個別に配置する。P1の\path{/etc/texnote}と\path{/etc/lilypondnote}は初期状態では存在
しないため、先に\texttt{install -d}で作る。

\begin{command}
# 管理用MacのTeXNoteリポジトリ直下で実行
scp Server/P1/p1.env mat@192.168.30.10:/tmp/texnote-p1.env

# P1で実行：ディレクトリ作成後、所有者root・600で配置
sudo install -d -o root -g root -m 750 /etc/texnote
sudo install -o root -g root -m 600 \
  /tmp/texnote-p1.env /etc/texnote/p1.env

# 管理用MacのLilyPondNoteリポジトリ直下で実行
scp Server/P1/p1.env mat@192.168.30.10:/tmp/lilypondnote-p1.env

# P1で実行
sudo install -d -o root -g root -m 750 /etc/lilypondnote
sudo install -o root -g root -m 600 \
  /tmp/lilypondnote-p1.env /etc/lilypondnote/p1.env
\end{command}

\texttt{install}は、ファイルをコピーすると同時に所有者と権限を設定する。秘密値を含むため、
\texttt{p1.env}を一般利用者から読める権限にしてはならない。

\section{systemdで常時起動する}

serviceファイルもMac側の各リポジトリに完成形で置いてある。抜粋を手入力せず、ファイル
全体を転送する。TeXとLilyPondではWorkingDirectory、EnvironmentFile、ExecStart、ポートが
異なる。

\lstinputlisting[language={},caption={TeXNoteの完全なtexnote-p1.service}]{/Users/mat/Documents/Git/TeXNote/Server/P1/texnote-p1.service}

\lstinputlisting[language={},caption={LilyPondNoteの完全なlilypondnote-p1.service}]{/Users/mat/Documents/Git/LilyPondNote/Server/P1/lilypondnote-p1.service}

\begin{command}
# 管理用MacのTeXNoteリポジトリ直下で実行
scp Server/P1/texnote-p1.service mat@192.168.30.10:/tmp/

# P1で実行
sudo install -o root -g root -m 644 /tmp/texnote-p1.service \
  /etc/systemd/system/texnote-p1.service

# 管理用MacのLilyPondNoteリポジトリ直下で実行
scp Server/P1/lilypondnote-p1.service mat@192.168.30.10:/tmp/

# P1で実行
sudo install -o root -g root -m 644 /tmp/lilypondnote-p1.service \
  /etc/systemd/system/lilypondnote-p1.service
\end{command}

\begin{command}
sudo systemctl daemon-reload
sudo systemctl enable --now texnote-p1.service lilypondnote-p1.service
systemctl status texnote-p1.service --no-pager -l
systemctl status lilypondnote-p1.service --no-pager -l
sudo ss -ltnp | grep -E ':8000|:8001'
\end{command}

\section{UFWとネットワーク境界}

P0からP1の二つの認証ポートだけを許可する。P1からP2へは、ルーターACLで8000、8001、
8002だけを許可する。

\begin{command}
sudo ufw default deny incoming
sudo ufw default allow outgoing
sudo ufw allow from 192.168.99.0/24 to any port 22 proto tcp
sudo ufw allow from 192.168.20.10 to any port 8000 proto tcp
sudo ufw allow from 192.168.20.10 to any port 8001 proto tcp
sudo ufw enable
sudo ufw status verbose
\end{command}

\section{初回疎通確認}

P0のLAN内入口をTCP 8080（TeXNote）とTCP 8081（LilyPondNote）に分けた場合の
ログイン確認例を次に示す。\texttt{192.168.20.10}は実際のP0 IPへ置き換える。

\begin{command}
curl -i -X POST http://192.168.20.10:8080/login \
  -H 'Content-Type: application/json' \
  -d '{"email":"tex-user@example.com","password":"REPLACE_ME","service":"texnote"}'

curl -i -X POST http://192.168.20.10:8081/login \
  -H 'Content-Type: application/json' \
  -d '{"email":"music-user@example.com","password":"REPLACE_ME","service":"lilypondnote"}'
\end{command}

正しい組合せはHTTP 200となる。次の誤った組合せは、仮に同じメールアドレスとパスワードが
両DBへ登録されていてもHTTP 422となる。

\begin{command}
curl -i -X POST http://192.168.20.10:8081/login \
  -H 'Content-Type: application/json' \
  -d '{"email":"tex-user@example.com","password":"REPLACE_ME","service":"texnote"}'
\end{command}

\begin{enumerate}
  \item TeXアカウントでTeXの\texttt{/login}と\texttt{/compile}を確認する。
  \item LilyPondアカウントで一度ログインし、同じJWTで\texttt{/compile}と
        \texttt{/transpose}を確認する。
  \item \texttt{/compile}がPDF、\texttt{/transpose}が移調後テキストだけを返すことを確認する。
  \item \texttt{/me}の\texttt{usageByService}に両方の回数が表示されることを確認する。
  \item DBを直接確認し、同じ\texttt{usage}テーブルにサービス種別が記録されることを確認する。
\end{enumerate}

\begin{command}
sudo -u lilypondnote-p1 sqlite3 \
  /opt/lilypondnote/P1/lilypondnote.db \
  'SELECT user_id, created_at, service, success FROM usage ORDER BY id DESC LIMIT 20;'
\end{command}

試験用利用者の登録方法は、次章ではなく文書末尾の「日常運用手続き」を参照する。

\section{完了チェックリスト}

\begin{itemize}
  \item[$\square$] TeXとLilyPondが別OSユーザー、別venv、別DBである。
  \item[$\square$] 各venvへ対応するrequirementsだけを導入した。
  \item[$\square$] P1にTeX LiveやLilyPond本体を導入していない。
  \item[$\square$] LilyPondの完成済み\texttt{main.py}に最初から\texttt{/transpose}がある。
  \item[$\square$] LilyPond PDFとリモート移調は同じDB・利用者・JWT・契約プラン・月間上限を共有する。
  \item[$\square$] ローカル移調はP1へ到達せず、リモート移調だけが利用回数へ記録される。
  \item[$\square$] LilyPondの\texttt{plans.py}にfree/standard/proの上限が定義されている。
  \item[$\square$] \texttt{usage.service}で\texttt{typeset}と\texttt{transpose}を判別できる。
  \item[$\square$] P0からだけP1の8000と8001へ接続できる。
  \item[$\square$] P1再起動後に両サービスが自動復旧する。
  \item[$\square$] 二つのDBをそれぞれ復元できるバックアップがある。
\end{itemize}

\clearpage
\section{日常運用手続き}

この章を、初回の試験利用者作成と、構築後の日常運用の両方に使用する。以下の操作は
\textbf{TeX DBとLilyPond DBのどちらを対象にしているか}を確認してから実行する。

\subsection{運用対象の早見表}

\begin{longtable}{@{}p{35mm}p{54mm}p{54mm}@{}}
\toprule
操作 & TeXNote & LilyPondNote\\
\midrule
DB & \path{/opt/texnote/P1/tex_service.db} & \path{/opt/lilypondnote/P1/lilypondnote.db}\\
登録 & \texttt{tools.useradd} & \texttt{useradd.py}\\
パスワード変更 & \texttt{tools.reset\_password} & \texttt{reset\_password.py}\\
契約値 & \texttt{plan}: free/standard/pro & \texttt{plan}: free/standard/pro\\
利用履歴 & TeXコンパイルだけ & PDFとリモート移調を\texttt{service}で区別\\
サービス & \texttt{texnote-p1.service} & \texttt{lilypondnote-p1.service}\\
\bottomrule
\end{longtable}

\subsection{DBをバックアップする}

SQLiteの\texttt{.backup}を使うため、通常はサービスを停止せず整合したバックアップを
取得できる。バックアップ先を最初に作る。

\begin{command}
sudo install -d -o texnote-p1 -g texnote-p1 -m 700 \
  /opt/texnote/P1/backups
sudo install -d -o lilypondnote-p1 -g lilypondnote-p1 -m 700 \
  /opt/lilypondnote/P1/backups
\end{command}

TeX DBをバックアップする。

\begin{command}
stamp=$(date +%Y%m%d-%H%M%S)
sudo -u texnote-p1 sqlite3 /opt/texnote/P1/tex_service.db \
  ".backup '/opt/texnote/P1/backups/tex_service-$stamp.db'"
sudo -u texnote-p1 sqlite3 \
  "/opt/texnote/P1/backups/tex_service-$stamp.db" 'PRAGMA integrity_check;'
\end{command}

LilyPond DBをバックアップする。PDF履歴と移調履歴は同じDBに入っているため、これ1個で
両方を保全できる。

\begin{command}
stamp=$(date +%Y%m%d-%H%M%S)
sudo -u lilypondnote-p1 sqlite3 /opt/lilypondnote/P1/lilypondnote.db \
  ".backup '/opt/lilypondnote/P1/backups/lilypondnote-$stamp.db'"
sudo -u lilypondnote-p1 sqlite3 \
  "/opt/lilypondnote/P1/backups/lilypondnote-$stamp.db" \
  'PRAGMA integrity_check;'
\end{command}

結果が\texttt{ok}であることを確認する。バックアップはP1内だけでなく、管理端末または
暗号化した別媒体へ定期的にコピーする。秘密情報を含むためGitへ登録しない。

\subsection{バックアップから復元する}

復元はDBを上書きするため、対象サービスを止め、現在のDBも退避してから行う。

\begin{command}
# LilyPond DBの復元例。TeXの場合はサービス名とパスをTeX用へ置き換える
sudo systemctl stop lilypondnote-p1.service
sudo cp -p /opt/lilypondnote/P1/lilypondnote.db \
  /opt/lilypondnote/P1/lilypondnote.db.before-restore
sudo cp -p /opt/lilypondnote/P1/backups/<復元するファイル>.db \
  /opt/lilypondnote/P1/lilypondnote.db
sudo chown lilypondnote-p1:lilypondnote-p1 \
  /opt/lilypondnote/P1/lilypondnote.db
sudo -u lilypondnote-p1 sqlite3 \
  /opt/lilypondnote/P1/lilypondnote.db 'PRAGMA integrity_check;'
sudo systemctl start lilypondnote-p1.service
systemctl status lilypondnote-p1.service --no-pager -l
\end{command}

\subsection{ユーザーを一覧・個別確認する}

TeX利用者を個別確認する。パスワードハッシュは表示しない。

\begin{command}
sudo -u texnote-p1 sh -c \
  'cd /opt/texnote/P1 && .venv/bin/python -m tools.check02 user@example.com'
\end{command}

TeX利用者を一覧表示する。

\begin{command}
sudo -u texnote-p1 sqlite3 -header -column \
  /opt/texnote/P1/tex_service.db \
  'SELECT id,email,status,plan,created_at FROM users ORDER BY id;'
\end{command}

LilyPond利用者と当月利用内訳を一覧表示する。

\begin{command}
sudo -u lilypondnote-p1 sqlite3 -header -column \
  /opt/lilypondnote/P1/lilypondnote.db \
  "SELECT u.id,u.email,u.status,u.plan,
          SUM(CASE WHEN h.service='typeset' THEN 1 ELSE 0 END) AS typeset_count,
          SUM(CASE WHEN h.service='transpose' THEN 1 ELSE 0 END) AS transpose_count
   FROM users AS u
   LEFT JOIN usage AS h ON h.user_id=u.id
     AND strftime('%Y-%m',h.created_at)=strftime('%Y-%m','now')
   GROUP BY u.id ORDER BY u.id;"
\end{command}

\lstinputlisting[caption={TeX利用者個別確認 tools/check02.py}]{/Users/mat/Documents/Git/TeXNote/Server/P1/progs/tools/check02.py}

\subsection{ユーザーを登録する}

TeX利用者は登録時にプランを選ぶ。入力できるプランは\texttt{free}、
\texttt{standard}、\texttt{pro}である。

\begin{command}
sudo -u texnote-p1 sh -c \
  'cd /opt/texnote/P1 && .venv/bin/python -m tools.useradd'
\end{command}

LilyPond利用者も登録時に\texttt{free}、\texttt{standard}、\texttt{pro}からプランを
選ぶ。何も入力しなければ\texttt{free}となり、月間上限は50回である。未定義のプラン名は
登録前に拒否される。

\begin{command}
sudo -u lilypondnote-p1 sh -c \
  'cd /opt/lilypondnote/P1 && .venv/bin/python useradd.py'
\end{command}

登録直後に一覧または個別照会で、メールアドレス、状態、契約値を確認する。サービス再起動は
不要である。

\lstinputlisting[caption={TeX利用者登録 tools/useradd.py}]{/Users/mat/Documents/Git/TeXNote/Server/P1/progs/tools/useradd.py}
\lstinputlisting[caption={LilyPond利用者登録 useradd.py}]{/Users/mat/Documents/Git/LilyPondNote/Server/P1/useradd.py}

\subsection{登録済みユーザーのパスワードを変更する}

これは管理者によるパスワード再設定である。利用者から本人確認を行った後、対象のメール
アドレスを一覧または個別照会で確認してから実行する。コマンド行へパスワードを書かず、
画面に表示されない対話入力で新しいパスワードを2回入力する。

TeX利用者のパスワードを変更する。

\begin{command}
sudo -u texnote-p1 sh -c \
  'cd /opt/texnote/P1 && .venv/bin/python -m tools.reset_password \
   user@example.com'
\end{command}

LilyPond利用者のパスワードを変更する。

\begin{command}
sudo -u lilypondnote-p1 sh -c \
  'cd /opt/lilypondnote/P1 && .venv/bin/python reset_password.py \
   user@example.com'
\end{command}

どちらも平文パスワードを保存せず、DBの\texttt{password\_hash}だけを新しいArgon2ハッシュ
へ置き換える。ユーザーの契約、状態、利用履歴は変えない。サービス再起動は不要である。
変更後は新しいパスワードでログインでき、古いパスワードでログインできないことを確認する。
既に発行済みのJWTはパスワード変更だけでは失効しないため、漏えい事故の場合はJWT署名鍵の
ローテーションも別途検討する。

\lstinputlisting[caption={TeXパスワード再設定 tools/reset\_password.py}]{/Users/mat/Documents/Git/TeXNote/Server/P1/progs/tools/reset_password.py}
\lstinputlisting[caption={LilyPondパスワード再設定 reset\_password.py}]{/Users/mat/Documents/Git/LilyPondNote/Server/P1/reset_password.py}

\subsection{利用停止と再開}

通常の解約・一時停止では履歴を残すため、ユーザー行を削除せず
\texttt{status='inactive'}へ変更する。停止後は新規ログインが拒否され、既発行JWTもAPI利用時
のDB確認で拒否される。

\begin{command}
# TeX利用者を停止
sudo -u texnote-p1 sqlite3 /opt/texnote/P1/tex_service.db \
  "UPDATE users SET status='inactive' WHERE email='user@example.com';"

# LilyPond利用者を停止
sudo -u lilypondnote-p1 sqlite3 /opt/lilypondnote/P1/lilypondnote.db \
  "UPDATE users SET status='inactive' WHERE email='user@example.com';"
\end{command}

再開時は対象DBで\texttt{status='active'}へ戻す。更新後は必ずSELECTまたは一覧表示で
対象者と状態を確認する。サービス再起動は不要である。

\subsection{ユーザーを物理削除する}

物理削除は監査・利用履歴も失うため、通常は行わない。削除が必要な場合は、直前バックアップ、
利用停止、対象メールアドレスの再確認を行う。

TeX DBは外部キーの連鎖削除を有効にして、関連履歴とAPIキーを削除する。

\begin{command}
sudo -u texnote-p1 sqlite3 /opt/texnote/P1/tex_service.db \
  "PRAGMA foreign_keys=ON;
   BEGIN IMMEDIATE;
   DELETE FROM users WHERE email='user@example.com' AND status='inactive';
   COMMIT;"
\end{command}

LilyPond DBは履歴を先に削除してからユーザーを削除する。

\begin{command}
sudo -u lilypondnote-p1 sqlite3 /opt/lilypondnote/P1/lilypondnote.db \
  "BEGIN IMMEDIATE;
   DELETE FROM usage WHERE user_id=(SELECT id FROM users
     WHERE email='user@example.com' AND status='inactive');
   DELETE FROM users WHERE email='user@example.com' AND status='inactive';
   COMMIT;"
\end{command}

削除後に一覧で不存在を確認する。誤削除した場合は、以後の書込みを止めてバックアップから
復元する。

\subsection{契約内容を更新する}

TeXとLilyPondのどちらも、回数そのものではなくプラン名を変更する。各上限は、それぞれの
サービスディレクトリに配置した\texttt{plans.py}に定義される。

\begin{command}
sudo -u texnote-p1 sh -c \
  'cd /opt/texnote/P1 && .venv/bin/python -m tools.change_plan \
   user@example.com standard'
sudo -u texnote-p1 sh -c \
  'cd /opt/texnote/P1 && .venv/bin/python -m tools.check02 user@example.com'
\end{command}

LilyPond利用者のプランも専用コマンドで変更する。

\begin{command}
sudo -u lilypondnote-p1 sh -c \
  'cd /opt/lilypondnote/P1 && .venv/bin/python change_plan.py \
   user@example.com standard'
sudo -u lilypondnote-p1 sqlite3 -header -column \
  /opt/lilypondnote/P1/lilypondnote.db \
  "SELECT id,email,status,plan FROM users
   WHERE email='user@example.com';"
\end{command}

変更は次のAPI要求から反映され、サービス再起動は不要である。LilyPondのPDFとリモート移調へ
別々の上限を設定する構造ではない。ローカル移調はP1を利用しないため、どのプランでも
カウント対象外である。

\lstinputlisting[caption={TeX契約プラン変更 tools/change\_plan.py}]{/Users/mat/Documents/Git/TeXNote/Server/P1/progs/tools/change_plan.py}
\lstinputlisting[caption={LilyPond契約プラン変更 change\_plan.py}]{/Users/mat/Documents/Git/LilyPondNote/Server/P1/change_plan.py}

\subsection{日常のサービス・ログ確認}

\begin{command}
systemctl is-active texnote-p1.service lilypondnote-p1.service
systemctl is-enabled texnote-p1.service lilypondnote-p1.service
journalctl -u texnote-p1.service -n 100 --no-pager
journalctl -u lilypondnote-p1.service -n 100 --no-pager
sudo ufw status verbose
df -h /opt/texnote/P1 /opt/lilypondnote/P1
\end{command}

ログへパスワード、JWT、内部トークン、楽譜本文を出力しない。異常時は、DB操作を繰り返す
前にログとバックアップを確保する。

\clearpage
\appendix
\section{旧サービスの停止・撤去}

この付録はP1の新規構築手順ではなく、以前のsystemdサービスやNginxが残っている場合の
調査・停止・撤去手順である。名前だけから推測して削除せず、必ず
\textbf{調査、内容確認、停止、自動起動解除、定義削除}の順で進める。現在必要な
\texttt{texnote-p1.service}と\texttt{lilypondnote-p1.service}を誤って停止・削除しない。

\subsection{以前のsystemdサービスを調べる}

動作中、登録済み、異常終了のサービスをそれぞれ表示する。

\begin{command}
systemctl list-units --type=service --state=running
systemctl list-unit-files --type=service
systemctl --failed

# 関係しそうな名前を絞り込む
systemctl list-unit-files --type=service | \
  grep -Ei 'tex|lily|python|uvicorn|fastapi|nginx'

# systemd以外から起動されたプロセスと待受ポートも確認する
ps aux | grep -Ei 'uvicorn|python|texnote|lilypond'
sudo ss -ltnp
\end{command}

候補を見つけたら、以下の\texttt{OLDUNIT}を実際のユニット名へ置き換える。
\texttt{systemctl cat}で起動コマンドを、\texttt{FragmentPath}でserviceファイルの実体を
確認する。確認が終わるまで削除してはならない。

\begin{command}
OLDUNIT=old-texnote.service
systemctl status "$OLDUNIT" --no-pager -l
systemctl cat "$OLDUNIT"
systemctl show "$OLDUNIT" -p FragmentPath -p DropInPaths
journalctl -u "$OLDUNIT" -n 100 --no-pager
\end{command}

\subsection{以前のsystemdサービスを止めて自動起動を解除する}

対象名と内容が間違いないことを確認した後、停止と自動起動解除を同時に行う。

\begin{command}
OLDUNIT=old-texnote.service
sudo systemctl disable --now "$OLDUNIT"
systemctl is-active "$OLDUNIT"
systemctl is-enabled "$OLDUNIT"
\end{command}

通常、確認結果が\texttt{inactive}と\texttt{disabled}なら停止・解除済みである。
待受ポートが残る場合は、別ユニット、cron、手動起動プロセス、コンテナ等がないかを
\texttt{sudo ss -ltnp}とプロセス一覧で再確認する。

\subsection{以前のsystemd定義を削除する}

\texttt{FragmentPath}が\path{/etc/systemd/system/old-texnote.service}のように
\path{/etc/systemd/system}配下であり、自分で作った旧定義だと確認できた場合だけ、その
表示されたパスを指定して削除する。次は例であり、実際のパスへ置き換える。

\begin{command}
sudo rm /etc/systemd/system/old-texnote.service
sudo systemctl daemon-reload
sudo systemctl reset-failed
systemctl status old-texnote.service
\end{command}

最後が\texttt{Unit ... could not be found}なら定義もなくなっている。
\path{/lib/systemd/system}または\path{/usr/lib/systemd/system}配下はパッケージ管理下である
可能性が高い。直接削除せず、所有パッケージを確認する。

\begin{command}
dpkg -S /lib/systemd/system/<対象>.service
\end{command}

パッケージが他の用途に不要であることを確認してから\texttt{apt remove}を検討する。
\texttt{DropInPaths}に追加設定が表示された場合も、内容と必要性を個別に確認してから扱う。

\section{以前のNginxの停止・撤去}

\subsection{Nginxの稼働状態と設定を調べる}

まずsystemd、プロセス、80/443番ポートの三方向から状態を確認する。

\begin{command}
systemctl status nginx --no-pager -l
systemctl is-enabled nginx
systemctl cat nginx
pgrep -a nginx
sudo ss -ltnp '( sport = :80 or sport = :443 or sport = :8000 )'
\end{command}

\texttt{systemctl cat nginx}で表示されるのはNginxプログラムの起動方法であり、
\texttt{server\_name}や\texttt{proxy\_pass}等のサイト設定ではない。次に設定ディレクトリと
有効なサイトを調べる。

\begin{command}
sudo find /etc/nginx -maxdepth 3 \
  \( -type f -o -type l \) -print
ls -l /etc/nginx/sites-enabled/
ls -l /etc/nginx/sites-available/
\end{command}

主な確認対象は\path{/etc/nginx/nginx.conf}、\path{/etc/nginx/sites-available}、
\path{/etc/nginx/sites-enabled}、\path{/etc/nginx/conf.d}である。Debian系の標準構成では、
\path{sites-available}は設定原本・候補の置場であり、\path{sites-enabled}からリンクされた
ものが有効になる。ただし、最終的に何が読み込まれるかは\texttt{nginx.conf}の
\texttt{include}と\texttt{nginx -T}で確定する。

\subsection{有効なサイト設定とリンク先を確認する}

例えば\path{/etc/nginx/sites-enabled/texnote-p1}がある場合、名前だけで旧設定と判断せず、
リンク先と内容を確認する。

\begin{command}
ls -l /etc/nginx/sites-enabled/
readlink -f /etc/nginx/sites-enabled/texnote-p1
sudo sed -n '1,240p' /etc/nginx/sites-enabled/texnote-p1
\end{command}

特に\texttt{listen}、\texttt{server\_name}、\texttt{location}、
\texttt{proxy\_pass}を確認する。次の関係なら、Nginxが80番の入口となり、同じ機械の旧
Uvicornへ中継する構成である。

\begin{command}
listen 80;
proxy_pass http://127.0.0.1:8000;
\end{command}

\path{sites-available/default}や\path{sites-available/texnote-p1.old}は、存在するだけでは
有効とは限らない。\path{sites-enabled}からリンクされているか、ほかの設定から直接
\texttt{include}されているかを確認する。

\begin{command}
sudo grep -Rni 'texnote-p1.old' /etc/nginx
sudo grep -RniE 'include|proxy_pass|127\.0\.0\.1:8000|server_name' \
  /etc/nginx/nginx.conf /etc/nginx/sites-enabled \
  /etc/nginx/sites-available /etc/nginx/conf.d
\end{command}

現在版と\texttt{.old}の違いを把握する場合は、削除前に差分を確認する。

\begin{command}
sudo diff -u /etc/nginx/sites-available/texnote-p1.old \
  /etc/nginx/sites-available/texnote-p1
\end{command}

\subsection{Nginxが実際に読み込んだ完成設定を確認する}

\texttt{nginx -T}は構文検査とともに、include後の全設定を表示する。ファイルが置かれて
いるだけなのか、実際に読み込まれているのかを確定するために必ず確認する。

\begin{command}
sudo nginx -T
sudo nginx -T 2>&1 | \
  grep -nE 'configuration file|listen|server_name|proxy_pass'
\end{command}

出力には証明書のパス、ホスト名、内部アドレス等が含まれ得るため、保存・共有時の扱いに
注意する。

\subsection{旧P1構成かを判定する}

次をすべて確認できた場合、旧TeXNoteのNginx・認証API構成である可能性が高い。

\begin{itemize}
  \item 現在の機械がP1の\texttt{192.168.30.10}である。
  \item Nginxが80番で待ち受ける。
  \item 有効な\texttt{texnote-p1}設定が\texttt{127.0.0.1:8000}へproxyする。
  \item \texttt{tex-auth.service}のUvicornが\texttt{127.0.0.1:8000}で待ち受ける。
  \item 新構成ではP0がNginxを担当し、P1はNginxを持たない。
\end{itemize}

次でP1のアドレス、旧認証サービス、待受を突き合わせる。

\begin{command}
hostname
ip -br address
systemctl show tex-auth.service \
  -p MainPID -p ExecStart -p FragmentPath
systemctl cat tex-auth.service
sudo ss -ltnp '( sport = :80 or sport = :8000 )'
\end{command}

一つでも不明なら停止せず、設定と利用者をさらに確認する。この機械がP0なら、新構成でも
Nginxを使用するため、旧P1だと思って停止してはならない。

\subsection{Nginx全体を止めて自動起動を解除する}

上記調査で旧P1構成と確定し、ほかのWebサイトやリバースプロキシがNginxを使用していない
ことを確認してから実行する。まず復旧用に設定全体を退避する。

\begin{command}
sudo cp -a /etc/nginx /etc/nginx.backup-before-p1-stop
sudo systemctl disable --now nginx
systemctl is-active nginx
systemctl is-enabled nginx
pgrep -a nginx
sudo ss -ltnp '( sport = :80 or sport = :443 or sport = :8000 )'
\end{command}

\texttt{inactive}、\texttt{disabled}となり、NginxプロセスとNginxによる80/443番待受が
なければ停止できている。80/443番が残る場合は、表示されたPIDとプログラム名から別の
Webサーバーを調べる。

\subsection{古いサイト設定だけを無効化する}

別用途でNginxを使い続ける場合は、Nginx全体を止めず、古いサイト設定だけを無効化する。
\path{/etc/nginx/sites-enabled}の対象がシンボリックリンクであり、リンク先も間違いないことを
\texttt{ls -l}で確認してから、enabled側だけを削除する。

\begin{command}
ls -l /etc/nginx/sites-enabled/
sudo rm /etc/nginx/sites-enabled/<古い設定名>
sudo nginx -t
sudo systemctl reload nginx
\end{command}

\texttt{nginx -t}が失敗した場合はreloadせず、削除したリンクを復元して設定を修正する。
\path{/etc/nginx/sites-available}側の原本は残るため、必要なら再びリンクを作成できる。

\subsection{Nginxパッケージ自体を削除する}

Nginxを今後まったく使わない場合に限る。まず設定を退避し、停止・解除する。

\begin{command}
sudo cp -a /etc/nginx /etc/nginx.backup-before-removal
sudo systemctl disable --now nginx
sudo apt remove nginx
\end{command}

設定ファイルも削除する場合は\texttt{purge}を使うが、復元が必要になる可能性を考慮する。

\begin{command}
sudo apt purge nginx nginx-common
sudo apt autoremove
\end{command}

最後に\texttt{systemctl status nginx}、\texttt{pgrep -a nginx}、
\texttt{sudo ss -ltnp}で、ユニット、プロセス、待受ポートを再確認する。

\section{現在のUFW状態と設定の確認}

この付録では、ルールを変更する前に、UFWが有効か、既定ポリシーが何か、どの通信を
許可しているかを調べる。SSH接続中に不用意にルールを削除すると、切断後に再接続できなく
なる可能性があるため、まず確認と退避だけを行う。

\subsection{有効・無効と既定ポリシーを確認する}

簡易表示と詳細表示を順に確認する。

\begin{command}
sudo ufw status
sudo ufw status verbose
\end{command}

\texttt{Status: active}なら有効、\texttt{Status: inactive}なら無効である。詳細表示では、
次の項目を確認する。

\begin{longtable}{@{}p{35mm}p{108mm}@{}}
\toprule
表示項目 & 意味\\
\midrule
\texttt{Logging} & 遮断等をUFWログへ記録する設定\\
\texttt{deny (incoming)} & 外部からの接続を既定では拒否する\\
\texttt{allow (outgoing)} & サーバーから外部への接続を既定では許可する\\
\texttt{routed} & この機械を経由して転送される通信の既定動作\\
\bottomrule
\end{longtable}

P1では通常、受信は既定拒否とし、P0からP1への必要ポートと管理ネットワークからのSSHだけを
明示的に許可する。

\subsection{許可・拒否ルールを番号付きで確認する}

\begin{command}
sudo ufw status numbered
sudo ufw show added
\end{command}

\texttt{status numbered}は現在の有効ルールを番号付きで表示する。\texttt{show added}は、
UFWへ追加されたルールをコマンドに近い形で表示する。IPv4とIPv6のルールが別々に表示される
場合がある。

番号指定でルールを削除できるが、番号は削除のたびに詰め直される。対象を十分確認し、1件
削除するたびに一覧を再表示する。

\begin{command}
sudo ufw status numbered
sudo ufw delete <削除する番号>
sudo ufw status numbered
\end{command}

この操作は接続中のSSHやサービスへ影響し得る。確認だけが目的なら\texttt{delete}は実行
しない。

\subsection{低レベルのルールを確認する}

UFWが生成したnetfilterルールを含めて確認する場合は次を使う。出力は多いため、通常の点検は
\texttt{status verbose}と\texttt{status numbered}で十分である。

\begin{command}
sudo ufw show raw
\end{command}

\subsection{アプリケーション定義を確認する}

\begin{command}
sudo ufw app list
sudo ufw app info OpenSSH
sudo ufw app info 'Nginx Full'
\end{command}

アプリケーション定義は、その名称がどのポートを意味するかを示すだけである。一覧に
\texttt{Nginx Full}が存在しても、実際に許可されているとは限らない。現在有効な許可は必ず
\texttt{sudo ufw status numbered}で確認する。

\subsection{UFWのsystemd状態を確認する}

\begin{command}
systemctl status ufw.service --no-pager -l
systemctl is-enabled ufw.service
\end{command}

UFWは起動時にルールを設定した後、常駐プロセスを必要としないため、
\texttt{active (exited)}と表示されても異常ではない。実際の有効・無効は
\texttt{sudo ufw status verbose}も併せて判断する。

\subsection{旧構成の80番・8000番ルールを確認する}

旧Nginxや旧認証APIで、80番または8000番を広く許可したルールが残っていないか確認する。

\begin{command}
sudo ufw status numbered
sudo ss -ltnp '( sport = :80 or sport = :443 or sport = :8000 )'
\end{command}

例えば\texttt{80/tcp ALLOW IN Anywhere}や\texttt{8000/tcp ALLOW IN Anywhere}は、全送信元
からの接続を許可する。ルールがあることと、実際にプロセスが待ち受けていることは別なので、
UFW一覧と\texttt{ss}の両方を確認する。なお、\texttt{grep ':80'}では8000番も一致するため、
上記の\texttt{ss}フィルターを使ってポートを区別する。

\subsection{変更前にSSH許可と設定退避を確認する}

SSHで作業している場合、最初に現在の接続元とSSH許可ルールを確認する。

\begin{command}
echo "$SSH_CONNECTION"
sudo ufw status numbered
\end{command}

管理端末のネットワークから22/tcpへの許可がない状態でUFWを有効化・再設定してはならない。
変更前のUFW設定を退避する場合は、日付を付けた明示的な保存先を使う。

\begin{command}
sudo cp -a /etc/ufw /etc/ufw.backup-before-p1-change
sudo ufw status verbose
sudo ufw status numbered
\end{command}

ルール変更後もSSHセッションをすぐ切断せず、別の端末から新規SSH接続できることを確認して
から作業を終了する。

\end{document}
